\chapter{A complete worked example}
\label{ch:worked}

\section{The three-location project}

The committed toy project has three locations and reciprocal road links. Each
directed edge also carries a second mode, which makes modal reallocation and
terminal congestion nontrivial. The inputs are small enough to inspect by
hand but use the same loader and closure engine as the empirical adapters.

\begin{figure}[H]
\centering
\includegraphics[width=0.72\textwidth]{toy-network.pdf}
\caption{Extended welfare elasticities in the three-location example. Line
width and color represent the model-implied benefit of a one-percent primitive
road-cost improvement.}
\label{fig:toy-network}
\end{figure}

Start from a fresh clone:
\begin{lstlisting}[language=bash]
julia --project=. -e 'using Pkg; Pkg.instantiate()'
julia --project=. bin/tnw.jl validate examples/toy/config.toml
julia --project=. bin/tnw.jl decompose examples/toy/config.toml
\end{lstlisting}
No environment variable, downloaded input, or source edit is required.

\section{Configuration}

The project declares its inputs and model:
\begin{lstlisting}
schema_version = 1
name = "toy"

[input]
adapter = "generic_csv_v1"
nodes = "data/nodes.csv"
edge_modes = "data/edge_modes.csv"

[model]
spatial_specification = "economic_geography"
alpha = 0.10
beta = -0.30
sigma = 9.0
modal_specification = "choice_logsum"
eta = 1.099
route_curvature = "theorem"

[policy]
mode = "road"
unit = "both"
shock_fraction = 0.01
\end{lstlisting}
The complete file also declares edge and terminal congestion, diagnostics, and
sensitivity paths. Paths are resolved relative to the configuration file, so
the project can be moved without editing absolute directories.

\section{Validation}

\code{validate} loads the input, builds the welfare-only closure, and prints a
machine-readable summary. The baseline has:
\begin{itemize}
  \item three nodes and six directed road policies;
  \item a route spectral radius of approximately \(0.0663\);
  \item a full spatial condition number of approximately \(1.16\times10^3\);
  \item matching bilateral margins and directed-edge traffic within the
  configured tolerance.
\end{itemize}

These checks do not establish that the toy calibration is empirically
plausible. They establish that it is a well-defined input to the model.

\section{Directed results}

Table \ref{tab:toy-results} reports the main directed quantities. All six road
edges have the same Hulten traffic share, but their extended effects differ
because their equilibrium incidence and pass-through differ.

\begin{table}[H]
\centering
\caption{Directed toy welfare elasticities}
\label{tab:toy-results}
\begin{tabular}{lrrr}
\toprule
Edge & Traditional & Extended primitive & Extended/traditional \\
\midrule
\(B\to C\) & 0.010000 & 0.007678 & 0.768 \\
\(C\to B\) & 0.010000 & 0.007644 & 0.764 \\
\(A\to C\) & 0.010000 & 0.006986 & 0.699 \\
\(A\to B\) & 0.010000 & 0.006817 & 0.682 \\
\(C\to A\) & 0.010000 & 0.006736 & 0.674 \\
\(B\to A\) & 0.010000 & 0.006597 & 0.660 \\
\bottomrule
\end{tabular}
\end{table}

For a one-percent primitive cost reduction on \(B\to C\), the local welfare
gain is approximately
\[
  0.01\times0.007678=0.00007678
\]
in log welfare, or \(0.007678\) percent. The derivative and the percentage
effect happen to share the same displayed number only because the policy
change is one percent.

\section{Traditional versus extended}

\begin{figure}[H]
\centering
\includegraphics[width=0.72\textwidth]{toy-scatter.pdf}
\caption{Traditional and extended welfare elasticities for the physical links
in the toy project. The 45-degree line is the efficient traffic benchmark.}
\label{fig:toy-scatter}
\end{figure}

The points lie below the 45-degree line in this calibration. This does not
imply a general theorem that equilibrium adjustment lowers benefits.
Parameters and network position determine the sign. The useful observation is
that equal traffic does not imply equal extended welfare once the baseline is
distorted.

\section{Decomposition}

\begin{figure}[H]
\centering
\includegraphics[width=0.92\textwidth]{toy-decomposition.pdf}
\caption{Exact closure components for the directed toy policies. Components
are signed levels, not percentage shares.}
\label{fig:toy-decomposition}
\end{figure}

The maximum inverse-gap error is \(3.3\times10^{-16}\), the closure-ladder
error is zero at displayed precision, and the maximum analytical channel
reconstruction error is \(2.8\times10^{-14}\). These diagnostics should be
checked before interpreting the bars.

\section{A routing-specific example}

The Braess-style project has four locations, reciprocal parallel routes, and a
connector. It is not a Wardrop-paradox replication; it is a controlled
demonstration that flexible and fixed route derivatives differ.

\begin{figure}[H]
\centering
\includegraphics[width=0.72\textwidth]{braess-network.pdf}
\caption{Extended welfare elasticities in the Braess-style routing example.}
\label{fig:braess}
\end{figure}

Run:
\begin{lstlisting}[language=bash]
julia --project=. bin/tnw.jl decompose examples/braess/config.toml
\end{lstlisting}
The route wedge is nonzero on the parallel branches. In the corresponding
unique-route test fixture, \code{FR=F} within \(10^{-10}\), as the theory
requires.

\section{Adapting the example}

Create a project skeleton:
\begin{lstlisting}[language=bash]
julia --project=. bin/tnw.jl init /path/to/my-network
\end{lstlisting}
Replace the CSV rows with model-ready data, retain the schema, and edit the
TOML values. The first successful milestone should be \code{validate}, not a
map or ranking. Once validation passes, run \code{analyze}; use
\code{decompose} only when the fixed-route memory estimate is acceptable.
